\begin{studiumworked}{M16. Escalas de vaporización y residencia · Intermedio}
\textbf{Enunciado.} Una gota académica aislada tiene diámetro inicial 120 micrómetros y constante $K=8\times10^{-7}$ m$^2$/s en la ley $d^2=d_0^2-Kt$. La residencia media del modelo es 0.030 s. Halle la duración de vaporización y compare al duplicar diámetro.

\textbf{Datos.} Diámetro inicial: $1.200000\times10^{-4}\,\mathrm{m}$; Constante: $8.000000\times10^{-7}\,\mathrm{m^2/s}$; Residencia: $0.03\,\mathrm{s}$.

\textbf{Incógnitas.} Magnitudes solicitadas y cierre físico.

\textbf{Hipótesis.} Modelo académico con hipótesis del capítulo. Datos originales; no representan hardware ni ensayo.

\textbf{Modelo físico.} Conservación y relaciones constitutivas declaradas.
\[t_v=d_0^2/K,\quad Da_{vap}=\tau_r/t_v\]
\textbf{Desarrollo.}
\begin{enumerate}
\item Convertir diámetro da $1.20\times10^{-4}$ m. Igualar diámetro final a cero produce $t_v=1.440000\times10^{-8}/8.000000\times10^{-7}=0.018$ s.
\item El cociente residencia/vaporización es $1.666667$. Con diámetro doble, $t_v$ se multiplica por cuatro: $0.072\,\mathrm{s}$, mayor que la residencia media.
\item La comparación advierte de una posible limitación del modelo, pero no demuestra que cada gota tenga exactamente la residencia media ni que una nube se comporte como una gota aislada.
\end{enumerate}
\textbf{Resultados.} $t_v=0.018$ s; con diámetro doble $0.072$ s.

\textbf{Interpretación.} El resultado corresponde a las hipótesis indicadas; su interpretación requiere conservar fronteras y unidades.

\textbf{Comprobación.} La constante tiene dimensión m²/s y el tiempo es positivo; la dependencia cuadrática se recupera por cocientes.

\textbf{Límites.} Sin calentamiento inicial, interacción entre gotas ni variación del entorno.
\end{studiumworked}
